\documentclass[11pt,a4paper]{article}
\usepackage[T1]{fontenc}
\usepackage[utf8]{inputenc}
\usepackage[english]{babel}
\usepackage{amsmath,amssymb}
\usepackage[margin=2cm]{geometry}
\setlength{\parindent}{0pt}
\setlength{\emergencystretch}{3em}
\allowdisplaybreaks

\title{About \texttt{correspondences.tex}}
\author{}
\date{}
\begin{document}
\maketitle
This note explains what the test data in \texttt{correspondences.tex} is used for. That file is input for the engine in \texttt{main.c}; it is read by the test \texttt{test\_maps\_and\_relations.c} in this folder.

\begin{enumerate}
\item \textbf{Purpose.} The file holds 16 statements about correspondences and binary relations, each to be proved. The statements are data: the engine knows none of them in advance.

\item \textbf{What is being tested.} Whether a set of pairs is the graph of a map, as for $\{(x, y) \in \mathbb{N} \times \mathbb{Z} \mid x = y^{4}\}$; reflexive, symmetric and transitive relations; the composition of two relations.

\item \textbf{How the statements are written.} Every property is written out as a formula, not named: for instance ``$f$ is injective'' is $\forall a, b\ (f(a) = f(b) \implies a = b)$. The engine must work from these definitions.

\item \textbf{How to run it.} From the top folder of the repository:
\begin{verbatim}
cc -std=c11 -O2 -Wall -Wextra -o /tmp/test_maps_and_relations \
    tests/set_theory/test_maps_and_relations.c && /tmp/test_maps_and_relations
\end{verbatim}

\item \textbf{Expected result.} One \texttt{PASS} line for the file when every statement is proved with a written proof. Today the engine proves 14 of the 16 statements: the first six, a set of pairs of numbers given by a rule, and whether it is the graph of a map; statement 8, a relation between the parts of a set, unfolded into a statement of chapter 14 and decided by regions (paragraph 17r of \texttt{parser.c}); and statements 9 to 14 and 16, about arbitrary relations on arbitrary sets, written with sets of pairs, inverses, compositions and the diagonal, unfolded (17p, 17q) and proved with instances of the hypotheses (chapter 28 of \texttt{main.c}), each with a certificate. Statement 7 relates elements to parts of a set, $\mathcal{P}(S)$ taken as a set of objects, and statement 15 is about given relations between numbers written as lists of pairs: both are reported as not supported, so the test fails. The file states the goal.
\end{enumerate}
\end{document}
